\CGZoneHeader{example-image-c}{NEGATIVOS}
\CGTextSingleColumn{}{Texto descritivo sobre o setor...}

\CGSectorHeader{guide_cyan}{NEGATIVOS}
\begin{CGSectorRoutes}
    \CGRoute{Efeito Colateral}{VIIIc}{30m}{Fixa}{}[EFEITO COLATERAL 8c| 30m (11+2) Faixa branca Tem uma via inacabada (Impermeável) à direita.]
    \CGRoute{Guaraçaí do grande}{IXa}{30m}{Fixa}{}[\CGstar1 GUARAÇAÍ DO GRANDE 9a| 26m (11+2) ! precisa de manutenção !]
    \CGRoute{Ilusionista}{VIIIc}{30m}{Fixa}{}[\CGstar3 ILUSIONISTA 8c| 26m (10+2) Via negativa. Foi regrampeada por Julio Francês (CE) em 2026 e adicionadas mais proteções. Video beta no canal de YT do Thiago Hara.]
    \CGRoute{Ojuara}{IXa}{30m}{Fixa}{}[\CGstar1 OJUARA 9a| 30m (10+2) Variante à esquerda do Ilusionista.]
    \CGRoute{Sementinha do Mal}{V}{10m}{Fixa}{}[\CGstar1 SEMENTINHA DO MAL 5º| 10m (3+2) Via mais fácil e curta do setor. Via de quarto grau fácil geral com crux protegido de quinto. Parada dupla em chapas PinGo.]
    \CGRoute{Segue o plano}{VIIa}{10m}{Fixa}{}[SEGUE O PLANO 6º | 10m (3+2) Via em grampos P. Regletes em via vertical e barriga levemente negativa. Parada dupla em grampos P (muito distantes um do outro). Entre duas arvores Via de Almir Neves (?)]
    \CGRoute{Fortalezas do conforto}{VIIb}{20m}{Fixa}{}[FORTALEZAS DO CONFORTO 7a| 20m (7+2) Crux no final exposto. Está planejado manutenção para protejer mais. Via de Almir Neves (?)]
    \CGRoute{Lombroso}{VIIc}{30m}{Fixa}{}[\CGstar1 LOMBROSO 7c| 30m (10+2)]
    \CGRoute{Tião Gavião}{VIIIa}{25m}{Fixa}{}[\CGstar1 TIÃO GAVIÃO 8a| 20m (6) Via com grau geral quinto e crux concentrado na barriga. Queda em platô, mas bem protegida. Via toda em grampos P. Rapel em proteção única.]
    \CGRoute{Coronel e o Pescador}{5° VIIa E2 D1}{110m}{Fixa}{}[\CGstar3 CORONEL E O PESCADOR 5º VIIa E2 D1 (110m) Ary Pacheco Neto
Equipamentos: 8 costuras, 1 corda de 70m Via mais classica do setor. Incluida por
Flávio Daflon no seu guia de 50 clássicas do país, sendo a única de Caiada.]
    \CGRoute{Penélope Charmosa}{VIIc}{30m}{Fixa}{}[\CGstar3 PENÉLOPE CHARMOSA 7c| 30m (9+2) Ary Pacheco [segunda enfiada inacabada]]
    \CGRoute{Dick Vigarista}{VIIc}{30m}{Fixa}{}[\CGstar3 DICK VIGARISTA 7c| 30m (10+2) Tem parada dupla intermediaria antes do primeiro crux (aprx. metade da via) em chapas PinGo. Originalmente graduada 7b, regraduada 7c/8a pelo conquistador em 2026. Júlio Francês (CE)]
    \CGRoute{Quinto Elemento}{VIIa}{}{Fixa}{}[\CGstar1 QUINTO ELEMENTO 7a| 30m (9+2) Esticado para pessoas pequenas Atenção para lance exposto após platô.]
    \CGRoute{A hora da verdade}{VIIb}{30m}{Fixa}{}[A HORA DA VERDADE 7b| 30m (8) Está em manutenção para troca de chapas (2026)]
    \CGRoute{Supertramp}{VIIIa}{30m}{Fixa}{}[SUPERTRAMP 7c/8a 30m (9+2) Manutenção iniciada em 2026. Adicionada uma chapa na saida. Ré-chapeada a primeira segunda e quarta chapa após a caverna/crux Almir 
seg Josias "O nome correto da via é Supertramp. Não tem nada a ver com Donald Trump. O nome vem de Alexander Supertramp, codinome do Chris McCandless, do Na Natureza Selvagem."]
    \CGRoute{Segundas Intenções}{VIIb}{}{Fixa}{}[\CGstar2 SEGUNDAS INTENÇÕES 7b| 30m (10+2) via bem protegida. Inicio de quinto grau com agarras, após a barriga segue lances de VII díficeis com boas agarras de mão. Parada dupla em grampos P | Rapel em diagonal desviando da vegetação.]
    \CGRoute{Pau de Rato}{V}{20m}{Fixa}{}[\CGstar1 PAU DE RATO 5º| 20m (5+2) Via de quinto grau, com crux no final difícil para o grau. Agarras parcialmente escondidas pela vegetação. Parada dupla em chapas PinGo. 
Precisa de manutenção, chapa da esquerda rodando.]
    \CGRoute{Desvio de Conduta}{5° VI E2}{110m}{Fixa}{}[\CGstar1 DESVIO DE CONDUTA 5º VI E2 D1 (105m) Ary Pacheco Neto, Clodoaldo Dante, Guilherme Almeida e Édervan Menger Equipamentos: 10 costuras, 2 cordas de 50 para rapel Acesso ao cume, ou rapel com 2 cordas. Lance delicado exposto na primeira enfiada]
    \CGRoute{Naturalmente}{Vsup}{30m}{Fixa}{}[\CGstar3 NATURALMENTE 5sup| 25m (6+2) Vsup mais recomendado de Caiada Via vertical, super protegida e com quedas limpas. Via tecnica com regletes. Parada dupla em chapas PinGo. Descida por Rapel. [Manutenção recomendada para adiconar argolas na parada para facilitar iniciantes]]
    \CGRoute{À Bientôt}{VIIc}{30m}{Fixa}{}[\CGstar2 ABIANTOT 7b/c| 30m (7+2) Inicialmente graduada 7a. Regraduada posteriormente após algumas (muitas) agarras terem quebrado. Crux de resistência com poucos descansos em agarras pequenas. Algumas pessoas não costuram a penúltima chapa seguindo pela direita; queda limpa porém muito exposto. O nome correto é "À Bientôt" 7c de acordo com Julio Francês, 7b de acordo com Thiago Hara.]
\end{CGSectorRoutes}

\CGSectorHeader{guide_cyan}{CAVERNONA}
\begin{CGSectorRoutes}
    \CGRoute{A Onda}{V}{20m}{Fixa}{}[A ONDA 5º| (3+2) Via mais a direita da Cavernona Via com crux na saida, considerar o uso de Clip stick. Virada interessante em degrau no fim da via. Atenção aos marimbondos escondidos. Fica à esqueda da Abiantot e de uma via inacabada. Termina em parada dupla com uma chapa e um grampo P. obs: o desenho do guia de 2022 está errado, a linha das chapas de via segue mais a esquerda do que indicado.]
    \CGRoute{Sete Vidas}{IVsup}{20m}{Fixa}{}[SETE VIDAS 4sup/9a| (4+2) Primeira parada simples em um unico grampo P compartilhada com a Ladrão de chapas.]
    \CGRoute{Ladrão de Chapas}{Vsup}{20m}{Fixa}{}[LADRÃO DE CHAPAS 5sup/8c| (5+2)]
    \CGRoute{El Patron}{Vsup}{25m}{Fixa}{}[EL PATRON 5sup/8c| (5+2)]
    \CGRoute{Até Breve Alicate}{VI}{25m}{Fixa}{}[ATÉ BREVE ALICATE 6º/9?| (4+2)]
    \CGRoute{Na covardia não}{Vsup}{25m}{Fixa}{}[NA COVARDIA NÃO 5sup/8a| (5+2)]
    \CGRoute{Orelha seca}{Vsup}{20m}{Fixa}{}[ORELHA SECA 5sup/9?| (5+2)]
    \CGRoute{Se escalar não case}{Vsup}{30m}{Fixa}{}[SE ESCALAR NÃO CASE 5sup| (5+2) Atenção as abelhas da cavernona. Algumas lacas e pedras frágeis]
\end{CGSectorRoutes}

